\chapter{Motion Estimation}

\textbf{Course:} Multimedia Communications — Università di Padova \textbf{Lecturer:} Marco Cagnazzo

\bigskip\hrule\bigskip

\section{1. Introduction}

\subsection{Motion and Apparent Motion}

Two distinct concepts:

\begin{itemize}
\item \textbf{Physical motion:} actual displacement of objects in 3D space (measurable by accelerometers, gyroscopes)
\item \textbf{Apparent motion (optical flow):} 2D displacement of pixel intensities between consecutive frames
\end{itemize}

These are \textbf{not} equivalent:
\begin{itemize}
\item Rotating sphere with uniform texture → physical motion exists, \textbf{zero} apparent motion
\item Moving light source → apparent motion without any physical object motion
\end{itemize}

We focus exclusively on \textbf{optical flow} — a 2D vector field $\mathbf{v}(n,m) = (u(n,m), v(n,m))^T$ describing per-pixel displacement between frames.

\begin{center}\includegraphics[width=0.600\linewidth]{Pasted_image_20260522134312.png}\end{center}
\subsection{Applications}

\begin{center}
\begin{tabular}{|l|l|}
\hline
\textbf{Domain} & \textbf{Application} \\ \hline
Video compression & Temporal prediction, reduces inter-frame redundancy \\ \hline
Video analysis & Object tracking, action recognition \\ \hline
Video enhancement & Deblurring, super-resolution, frame interpolation \\ \hline
Medical imaging & Cardiac motion analysis \\ \hline
Autonomous driving & Obstacle detection \\ \hline
\end{tabular}
\end{center}

\bigskip\hrule\bigskip

\section{2. Variational Methods}

Also called \textbf{differential} or \textbf{gradient-based} methods. Produce a \textbf{dense} motion field — one vector per pixel. Key representative: \textit{Horn and Schunck}.

General formulation as energy minimization:

\[
\hat{v} = \arg\min_v E(v)
\]

\[
E(v) = E_{\text{data}}(v) + \lambda\, E_{\text{smooth}}(v)
\]
\subsection{Optical Flow Equation}

\textbf{Problem:} given frames $f_k$ (current) and $f_h$ (reference, $h = k \pm 1$), find $\mathbf{v}(n,m)$ such that:

\[
f_k(n,m) \approx f_h(n - u(n,m),\; m - v(n,m))
\]

\begin{calloutbox}{red!70!black}{Constant Illumination Hypothesis}
\[
f_k(n,m) = f_h(n - u,\; m - v)
\]
Pixel brightness does not change between frames. This is the foundational assumption of all gradient-based ME methods. It is violated by lighting changes, reflections, and occlusions.
\end{calloutbox}

\textbf{Derivation via Taylor expansion:}

\[
f_h(n-u, m-v) \approx f_h(n,m) - u\,f_n - v\,f_m
\]

Substituting into the constant illumination hypothesis:

\[
f_k - f_h = u\,f_n + v\,f_m \implies f_t = u\,f_n + v\,f_m
\]

\begin{calloutbox}{red!70!black}{Optical Flow (OF) Equation}
\[
\nabla f \cdot \mathbf{v} + f_t = 0
\]
One equation, two unknowns $(u, v)$ — \textbf{underdetermined}. This is the \textbf{aperture problem}: along an edge, only the normal component of motion is observable from local measurements.
\end{calloutbox}

\subsection{Horn and Schunck Algorithm}

Horn and Schunck introduced a constraint on the \textbf{total variation} of the velocity over a region $\mathcal{R}$. The original formulation is a constrained minimization:

\[
\iint_{\mathcal{R}} (u f_x + v f_y + f_t)^2 \, dx\, dy = \min
\]

\[
\iint_{\mathcal{R}} \left(\|\nabla u\|^2 + \|\nabla v\|^2\right) dx\, dy \leq \tau
\]

Via Lagrange multiplier, this becomes the unconstrained problem of minimizing $J$:

\[
J = \iint_{\mathcal{R}} (u f_x + v f_y + f_t)^2 \, dx\, dy + \lambda \left[\iint_{\mathcal{R}} \left(\|\nabla u\|^2 + \|\nabla v\|^2\right) dx\, dy - \tau\right]
\]

Equivalently, minimize (dropping the constant $-\lambda\tau$):

\begin{calloutbox}{red!70!black}{Horn and Schunck Energy Functional}
\[
E(u,v) = \iint_{\mathcal{R}} \left[(\nabla f \cdot \mathbf{v} + f_t)^2 + \lambda^2\left(|\nabla u|^2 + |\nabla v|^2\right)\right] dn\, dm
\]
\begin{itemize}
\item First term: data fidelity (OF equation residual)
\item Second term: smoothness regularization weighted by $\lambda^2$
\item Problem is \textbf{convex} → unique global minimum
\end{itemize}
\end{calloutbox}

\textbf{Euler-Lagrange equations} derived from the functional:

\[
f_n(f_n u + f_m v + f_t) = \lambda^2 \Delta u
\]
\[
f_m(f_n u + f_m v + f_t) = \lambda^2 \Delta v
\]

Since $\Delta u \approx \bar{u} - u$ (Laplacian via local average), one obtains the linear system:

\[
\begin{cases} (f_n^2 + \lambda)u + f_n f_m v = \lambda\bar{u} - f_n f_t \\ f_n f_m u + (f_m^2 + \lambda)v = \lambda\bar{v} - f_m f_t \end{cases}
\]

\textbf{Iterative solution (Gauss-Seidel):}

\[
u^{k+1} = \bar{u}^k - f_n \frac{f_n \bar{u}^k + f_m \bar{v}^k + f_t}{\lambda^2 + f_n^2 + f_m^2}
\]

\[
v^{k+1} = \bar{v}^k - f_m \frac{f_n \bar{u}^k + f_m \bar{v}^k + f_t}{\lambda^2 + f_n^2 + f_m^2}
\]

where $\bar{u}^k, \bar{v}^k$ are local spatial averages at iteration $k$.

\begin{center}
\begin{minipage}{0.330\linewidth}\centering\includegraphics[width=\linewidth]{Pasted_image_20260522134513.png}\end{minipage}\hfill\begin{minipage}{0.330\linewidth}\centering\includegraphics[width=\linewidth]{Pasted_image_20260522134536.png}\end{minipage}\hfill\begin{minipage}{0.330\linewidth}\centering\includegraphics[width=\linewidth]{Pasted_image_20260522134558.png}\end{minipage}
\end{center}

\textbf{Conclusions:}

\begin{center}
\begin{tabular}{|l|l|}
\hline
\textbf{\textbf{Advantages}} & \textbf{Dense field (1 vector/pixel), convex optimization (unique solution), good in low-texture areas} \\ \hline
\textbf{Disadvantages} & Over-smoothed at motion boundaries, slow for large displacements, brightness constancy violated in practice \\ \hline
\end{tabular}
\end{center}

\textbf{Applications:} object tracking, video compression, motion analysis. Foundational approach in computer vision — basis for more advanced and modern optical flow estimation algorithms.

\bigskip\hrule\bigskip

\section{3. Block-Matching Methods}

\textbf{Core idea:} divide image into rectangular blocks; for each block in current frame $f_k$, find best-matching block in reference frame $f_h$. Motion described by one translation vector per block → \textbf{piecewise-constant} \textit{Motion Vector Field (MVF)}.

Foundation of all major video compression standards (MPEG, H.264, H.265, AV1).

\subsection{Notation and Formulation}

\begin{itemize}
\item $N \times M$ image; blocks of size $P \times Q$
\item Block $B_{p,q} = \{p, p+1, \ldots, p+P-1\} \times \{q, q+1, \ldots, q+Q-1\}$
\item $\mathbf{f}_k(B_{p,q})$: vector of $PQ$ luminance values in $B_{p,q}$ at time $k$
\item Search window $\mathcal{W}$: set of candidate displacement vectors
\end{itemize}

\begin{calloutbox}{red!70!black}{Block-Matching Optimization}
\[
\left(\hat{i}, \hat{j}\right) = \arg\min_{(i,j)\in\mathcal{W}}\; d\!\left[\mathbf{f}_k(B_{p,q}),\; \mathbf{f}_h(B_{p-i,q-j})\right]
\]
BM methods differ in: choice of criterion $d$, set $\mathcal{W}$ (search window + strategy), block size and shape.
\end{calloutbox}

The full MVF is then:

\[
\forall (n,m) \in B_{p,q}: \quad (u(n,m), v(n,m)) = \arg\min_{(i,j)\in\mathcal{W}} d\!\left[\mathbf{f}_k(B_{p,q}),\; \mathbf{f}_h(B_{p-i,q-j})\right]
\]

A more precise notation uses subscript $h \to k$ to indicate reference frame $h$ and current frame $k$.
\begin{itemize}
\item \textbf{Forward motion} ($h = k-1$): predict from previous frame
\item \textbf{Backward motion} ($h = k+1$): predict from next frame
\end{itemize}
\begin{center}\includegraphics[width=0.600\linewidth]{Pasted_image_20260522134710.png}\end{center}
\subsection{Evaluation of the MVF}

Three complementary metrics:

\textbf{1. Quality — MC-ed MSE / PSNR}

Motion-compensated prediction of pixel $(n,m)$:

\[
\tilde{f}_k(n,m) = f_h\!\left(n + u_{h\to k}(n,m),\; m + v_{h\to k}(n,m)\right)
\]

Prediction error: $e(n,m) = f_k(n,m) - \tilde{f}_k(n,m)$

\[
\mathcal{E} = \frac{1}{NM}\sum_{n,m} e^2(n,m) \qquad \text{PSNR} = 10\log_{10}\frac{255^2}{\mathcal{E}} \quad \text{[dB]}
\]

\textbf{2. Coding cost — MVF entropy}

Bits to losslessly encode the MVF. Techniques: fixed-length, Exp-Golomb, Huffman, arithmetic coding. Technique-independent measure: \textbf{empirical entropy} of $(u,v)$ marginal distributions.

\textbf{3. Computational complexity}

To compute $\arg\min$ over $\mathcal{W}$, for each of $\lceil NM/PQ \rceil$ blocks, evaluate $|\mathcal{W}|$ instances of $d[\cdot,\cdot]$. Complexity drivers:
\begin{itemize}
\item Block size $P \times Q$ → number of blocks
\item $|\mathcal{W}|$ → search window size and strategy
\item Cost of computing $d$
\end{itemize}

\textbf{Block size trade-off:}

\begin{center}
\begin{tabular}{|l|l|l|l|}
\hline
\textbf{Block size} & \textbf{Complexity} & \textbf{Coding cost} & \textbf{MSE} \\ \hline
Large ($\uparrow$) & Low (fewer blocks) & Low (fewer vectors) & High (less flexible model) \\ \hline
Small ($\downarrow$) & High & High & Low \\ \hline
\end{tabular}
\end{center}

\begin{center}\fbox{\parbox{0.85\linewidth}{\centering\itshape Figure placeholder: MVF comparison: block sizes 32×32, 16×16, 8×8}}\end{center}

\subsection{Cost Functions: SAD and SSD}

\begin{calloutbox}{red!70!black}{Norm-Based Dissimilarity (General)}
\[
J(i,j) = \|\mathbf{f}_k(B_{p,q}) - \mathbf{f}_h(B_{p-i,q-j})\|_p^p \tag{2}
\]
Two principal cases: $p=1$ ($\mathcal{L}^1$ norm → SAD), $p=2$ ($\mathcal{L}^2$ norm → SSD)
\end{calloutbox}

\textbf{SSD} — Sum of Squared Differences ($p = 2$):

\[
J_{\text{SSD}}(i,j) = \sum_{(n,m)\in B_{p,q}} \left[f(n,m,k) - f(n-i,\,m-j,\,h)\right]^2
\]

\begin{itemize}
\item Minimizes energy of MC-ed prediction error → \textbf{best quality} metric
\item Complex to compute (requires multiplications)
\item Sensitive to \textbf{outliers}: a single pixel violating brightness constancy can corrupt the vector
\item Irregular MVF → higher entropy → higher coding cost
\end{itemize}

\textbf{SAD} — Sum of Absolute Differences ($p = 1$):

\[
J_{\text{SAD}}(i,j) = \sum_{(n,m)\in B_{p,q}} \left|f(n,m,k) - f(n-i,\,m-j,\,h)\right|
\]

\begin{itemize}
\item More robust to outliers (linear, not quadratic penalization)
\item Produces more \textbf{regular} MVF → lower coding cost
\item Slightly worse MSE than SSD (by construction: SAD minimizes $\mathcal{L}^1$, not $\mathcal{L}^2$ error)
\end{itemize}

\begin{calloutbox}{green!55!black}{SAD vs SSD on \textit{Flower Garden} sequence (16×16 blocks)}
\begin{center}
\begin{tabular}{|l|l|l|}
\hline
\textbf{Criterion} & \textbf{Rate (MV)} & \textbf{PSNR(Pred)} \\ \hline
SSD & 2143 bits & 22.46 dB \\ \hline
SAD & 2103 bits & 22.30 dB \\ \hline
\end{tabular}
\end{center}
SSD gives better PSNR but costs more bits; SAD gives more regular MVF.
\end{calloutbox}

\begin{center}\includegraphics[width=0.600\linewidth]{Pasted_image_20260522134839.png}\end{center}
\begin{center}\includegraphics[width=0.600\linewidth]{Pasted_image_20260522134856.png}\end{center}
\begin{center}\includegraphics[width=0.600\linewidth]{Pasted_image_20260522134904.png}\end{center}
\begin{center}\includegraphics[width=0.600\linewidth]{Pasted_image_20260522134913.png}\end{center}
\begin{center}\includegraphics[width=0.600\linewidth]{Pasted_image_20260522134931.png}\end{center}
\subsection{Regularized Cost Function}

Penalize vectors that differ too much from neighbors:

\[
J_{\text{REG}}(i,j) = \|\mathbf{f}_k(B_{p,q}) - \mathbf{f}_h(B_{p-i,q-j})\|_p^p + \lambda\, R(i,j)
\]

where $R(i,j)$ = distance of $(i,j)$ from neighborhood representative (average or median).

More generally, encode coding cost directly:

\[
J(v) = d(v) + \lambda_{ME}\, r(v)
\]

\begin{itemize}
\item $r(v)$: bits to encode vector $v$
\item $\lambda_{ME}$ controls rate-distortion trade-off:
\begin{itemize}
\item $\lambda_{ME} \to 0$: only minimize distortion $d$
\item $\lambda_{ME} \to \infty$: only minimize rate $r$
\end{itemize}
\end{itemize}

\begin{calloutbox}{green!55!black}{Regularization effect on \textit{Flower Garden}}
\begin{center}
\begin{tabular}{|l|l|l|}
\hline
\textbf{} & \textbf{Rate (MV)} & \textbf{PSNR(Pred)} \\ \hline
SSD (no reg.) & 2143 bits & 22.46 dB \\ \hline
SSD (regularized) & 2008 bits & 22.35 dB \\ \hline
\end{tabular}
\end{center}
Rate reduced 6.3\% at cost of 0.11 dB.
\end{calloutbox}

\begin{center}\includegraphics[width=0.600\linewidth]{Pasted_image_20260522135013.png}\end{center}
\subsection{Search Strategies}

\subsubsection{Full Search (FS)}

Search window:

\[
\mathcal{W} = \{-A,\ldots,A\} \times \{-B,\ldots,B\}
\]

Test all $n^2 = (2A+1)^2$ candidates. \textbf{Guarantees global minimum.} Very costly.

The cost function surface $J(i,j)$ is a 3D landscape:
\begin{itemize}
\item \textbf{Global minimum}: optimal MV
\item \textbf{Local minima}: traps for fast methods
\item \textbf{Flat regions}: ambiguous (low texture)
\end{itemize}
\subsubsection{Three Steps Search (3SS)}

Assumes \textbf{unimodal} cost surface. Steps with decreasing stride $D, D/2, \ldots, 1$; tests 4 or 8 points per step.
\begin{itemize}
\item \textasciitilde{}25 tests vs 225 for FS (-89\%)
\item Risk: trapped in local minimum if surface is multimodal
\end{itemize}

\subsubsection{Diamond Search (DS)}

Uses \textbf{Large Diamond Search Pattern (LDSP)} — moves until minimum is at center, then refines with \textbf{Small Diamond (SDSP)}.
\begin{itemize}
\item 3–5 new tests per iteration (9 for first step)
\item \textasciitilde{}23 tests (-90\% vs FS)
\item Unbounded: can follow large motions across frame
\end{itemize}
\begin{center}\includegraphics[width=0.600\linewidth]{Pasted_image_20260522135115.png}\end{center}
\begin{center}\includegraphics[width=0.600\linewidth]{Pasted_image_20260522135125.png}\end{center}
\subsubsection{Hexagon Search}

H.264/AVC standard. 6-sided \textbf{Large Hexagon Pattern (LHP)}; only 3 new points per step when minimum ≠ center. Final refinement with \textbf{Small Hexagon Pattern (SHP)}.
\begin{itemize}
\item \textasciitilde{}17 tests (-92\% vs FS)
\item More efficient than diamond for natural video
\end{itemize}
\begin{center}\includegraphics[width=0.600\linewidth]{Pasted_image_20260522135139.png}\end{center}
\begin{center}\includegraphics[width=0.600\linewidth]{Pasted_image_20260522135148.png}\end{center}
\subsubsection{Comparison}

\begin{center}
\begin{tabular}{|l|l|l|}
\hline
\textbf{Strategy} & \textbf{Tests} & \textbf{Δ\% vs FS} \\ \hline
Full Search (FS) & 225 & Baseline \\ \hline
2D-Log Search & 25 & \textasciitilde{}−89\% \\ \hline
Diamond Search & 23 & \textasciitilde{}−90\% \\ \hline
Hexagon Search & 17 & \textasciitilde{}−92\% \\ \hline
\end{tabular}
\end{center}

\textit{All converge to same optimal vector in this example.}

\textbf{Key distinctions:}
\begin{itemize}
\item FS: bounded search, global optimum guaranteed
\item Diamond/Hexagon: \textbf{iterative, unbounded} — can follow motion of any amplitude; risk of local minima (rare in natural video)
\end{itemize}

\subsubsection{Modern Fast Search: TZSearch}

Reference algorithm for recent standards (H.265, VVC), handles blocks up to 64×64.

Three phases:
\begin{enumerate}
\item \textbf{Search Predictors:} test vectors from spatial/temporal neighbors; stop if error is small
\item \textbf{Adaptive Loop:} if no predictor is good, apply Diamond/Square with increasing step size
\item \textit{(Full refinement phase — details not covered in slides)}
\end{enumerate}

\textbf{Key advantage:} more robust to local minima than pure Diamond/Hexagon.

\subsection{Subpixel Precision}

Motion is not constrained to integer pixel displacement. Half-pixel, quarter-pixel vectors improve prediction quality.

\textbf{Hierarchical approach:}
\begin{enumerate}
\item Find integer $(\hat{i}, \hat{j}) \in \mathbb{Z}^2$
\item Test $(\hat{i} \pm \tfrac{1}{2},\; \hat{j} \pm \tfrac{1}{2})$
\item Test $(\hat{i} \pm \tfrac{1}{4},\; \hat{j} \pm \tfrac{1}{4})$, etc.
\end{enumerate}

Non-integer positions evaluated via \textbf{bilinear interpolation}:

\[
f(n+a, m+b) = (1-a)(1-b)\,x + a(1-b)\,y + (1-a)b\,z + ab\,w
\]

where $x,y,z,w$ are the four integer-pixel neighbors. Higher-order filter interpolation also used.
\begin{center}\includegraphics[width=0.600\linewidth]{Pasted_image_20260522135238.png}\end{center}
\subsection{Variable Block Size}

\textbf{Problem:} homogeneous-motion assumption fails at object borders.

\begin{itemize}
\item \textbf{Solution 1 — fixed smaller blocks:} better precision, higher complexity and coding cost
\item \textbf{Solution 2 — variable-size blocks (adopted in modern standards):} split block only when needed
\end{itemize}

\textbf{Hierarchical splitting algorithm:}

\textbf{Input:} Frame $F$, reference $R$, initial block size $B$, $\lambda$ \textbf{Output:} MVF + block partitioning

\begin{verbatim}
for each block b of size B in F do
    Calculate J(v) = D + λR for b;
    Divide b into four sub-blocks;
    for each sub-block s_i do
        Calculate J_i(v) = D_i + λR_i for s_i;
    end
    J_sub = Σ^4_{i=1} J_i(v)
    if J_sub < J(v) then
        Apply algorithm recursively to each sub-block;
    else
        Keep b; store MV;
    end
end

\end{verbatim}

\begin{itemize}
\item $D$: distortion (SAD or SSD); $R$: rate (bits for MV); $\lambda$: Lagrange multiplier
\item Terminates when: minimum block size reached, or $J(v)$ no longer decreases
\item Balances accuracy vs. coding efficiency
\end{itemize}

\bigskip\hrule\bigskip

\section{4. Parametric Methods}

\textbf{Definition:} MVF modeled as a \textbf{closed-form function} of pixel position. Degrees of freedom = model parameters.

Scope:
\begin{itemize}
\item \textbf{Global model:} one set of parameters for the entire scene
\item \textbf{Per-region model:} region = object (requires segmentation) or block
\end{itemize}

\begin{calloutbox}{red!70!black}{Relationship to Block Matching}
BM is a special case of parametric ME: MVF is block-wise constant (pure translation), 2 parameters per block. More general models trade off:
\begin{itemize}
\item Fewer parameters → robust, but may miss motion complexity
\item More parameters → general, but prone to overfitting
\end{itemize}
\end{calloutbox}

\subsection{Translational Model}

Rigid object translating orthogonally to the optical axis, no rotation.

\[
\mathbf{v}(\mathbf{p}) = \begin{bmatrix} v_x \\ v_y \end{bmatrix} = \begin{bmatrix} b_1 \\ b_2 \end{bmatrix}
\]

Constant field — only 2 parameters. Same as block matching with $P = N$, $Q = M$.

\subsection{Affine Model}

\begin{calloutbox}{red!70!black}{Affine Motion Model}
\[
\mathbf{v}(\mathbf{p}) = \mathbf{b} + \mathbf{B}\mathbf{p} = \begin{bmatrix} b_1 \\ b_2 \end{bmatrix} + \begin{bmatrix} b_3 & b_4 \\ b_5 & b_6 \end{bmatrix} \mathbf{p}
\]
\begin{itemize}
\item 6 parameters: 2 for translation $\mathbf{b}$, 4 for linear part $\mathbf{B}$
\item Special case $\mathbf{B} = 0$: pure translation
\item Represents translation, zoom, rotation, shear in combination
\item Small parameter count → robust estimation
\end{itemize}
\end{calloutbox}

\textbf{Examples:}

\begin{center}
\begin{tabular}{|l|l|l|l|}
\hline
\textbf{Motion} & \textbf{$\mathbf{b}$} & \textbf{$\mathbf{B}$} & \textbf{Field} \\ \hline
Translation & $[0.5,\; 2]^T$ & $\mathbf{0}$ & Uniform $v_x=0.5$, $v_y=2$ \\ \hline
Zoom in & $\mathbf{0}$ & $\text{diag}(0.5, 0.5)$ & $v_x=0.5x$, $v_y=0.5y$ \\ \hline
Zoom out & $\mathbf{0}$ & $\text{diag}(-0.5,-0.5)$ & $v_x=-0.5x$, $v_y=-0.5y$ \\ \hline
Rotation & $\mathbf{0}$ & $\begin{bmatrix}0&-0.5\\0.5&0\end{bmatrix}$ & $v_x=-0.5y$, $v_y=0.5x$ \\ \hline
\end{tabular}
\end{center}

\begin{center}
\begin{minipage}{0.450\linewidth}\centering\includegraphics[width=\linewidth]{Pasted_image_20260522135333.png}\end{minipage}\hfill\begin{minipage}{0.450\linewidth}\centering\includegraphics[width=\linewidth]{Pasted_image_20260522135458.png}\end{minipage}
\end{center}

\begin{center}
\begin{minipage}{0.500\linewidth}\centering\includegraphics[width=\linewidth]{Pasted_image_20260522135546.png}\end{minipage}\hfill\begin{minipage}{0.500\linewidth}\centering\includegraphics[width=\linewidth]{Pasted_image_20260522135610.png}\end{minipage}
\end{center}
Only 6 parameters encode the entire field — replacing hundreds of block-matching vectors.

\subsection{Parameter Estimation}

\subsubsection{Indirect Estimation}

\begin{enumerate}
\item Estimate dense field $(u(n,m), v(n,m))$ first (e.g., with Horn-Schunck)
\item Fit model parameters via least squares:
\end{enumerate}

\[
\pi^* = \arg\min_\pi \sum_{n,m \in \mathcal{R}} \left[u - u_\pi\right]^2 + \left[v - v_\pi\right]^2
\]

\begin{itemize}
\item For affine: $\pi = [b_1\, b_2\, b_3\, b_4\, b_5\, b_6]$, solvable with gradient descent
\item \textbf{Disadvantage:} strongly dependent on quality of initial dense estimation; needs region $\mathcal{R}$ with homogeneous motion
\end{itemize}

\subsubsection{Direct Estimation}

Parameters enter the estimation directly. Two approaches:

\textbf{Via OF equation:}

\[
\pi^* = \arg\min_\pi \sum_{n,m \in \mathcal{R}} \left[u_\pi(n,m)\, f_x + v_\pi(n,m)\, f_y + f_t\right]^2
\]

\textbf{Via SSD/SAD on parametrized MVF:}

\[
\pi^* = \arg\min_\pi \sum_{n,m \in \mathcal{R}} \left[e(n,m)\right]^2
\]

\[
e(n,m) = f\!\left(n - u_\pi(n,m),\, m - v_\pi(n,m),\, t-1\right) - f(n,m,t)
\]

\textit{Note:} SAD-based block matching = direct parametric estimation with pure translation model ($\mathbf{B} = 0$).

\bigskip\hrule\bigskip

\section{5. Deep Learning Methods}

\subsection{The Paradigm Shift}

Traditional methods use hand-crafted models (gradient, block matching). DL treats ME as a \textbf{supervised or unsupervised learning problem} via CNNs.

\textbf{Key requirements:}
\begin{itemize}
\item Specialized architectures for pixel-level correspondence
\item Large-scale training data (synthetic or real)
\item GPU inference
\end{itemize}

\textbf{Ground truth challenge:} optical flow GT is hard to obtain for real sequences.

Solutions:
\begin{itemize}
\item \textbf{Synthetic datasets:} FlyingChairs, Sintel (perfect labels from rendering)
\item \textbf{Data augmentation:} geometric and photometric transforms
\item \textbf{Unsupervised:} minimize photometric error between current frame and warped reference
\end{itemize}

\subsection{FlowNet (2015)}

First end-to-end CNN for optical flow.
\begin{center}\includegraphics[width=0.600\linewidth]{Pasted_image_20260522135655.png}\end{center}

Two variants:
\begin{itemize}
\item \textbf{FlowNetS:} concatenate both frames → single encoder-decoder
\item \textbf{FlowNetC:} explicit \textbf{correlation layer} between feature maps of each frame
\end{itemize}

\textgreater{} Dosovitskiy et al. (2015). \textit{FlowNet: Learning Optical Flow with Convolutional Networks}

\textbf{Performance \& limits:}

\begin{center}
\begin{tabular}{|l|l|}
\hline
\textbf{Aspect} & \textbf{Detail} \\ \hline
Speed & 10–100 FPS on mid-range GPU \\ \hline
Weakness & Struggles with small displacements and fine details \\ \hline
Data dependency & Heavily tied to \textit{Flying Chairs} quality \\ \hline
Failure mode & Repetitive patterns (limited CNN receptive field) \\ \hline
\end{tabular}
\end{center}

\subsection{FlowNet 2.0 (2017)}

\textgreater{} Ilg et al. (2017). \textit{FlowNet 2.0: Evolution of Optical Flow Estimation with Deep Networks}

\textbf{Improvements:}
\begin{itemize}
\item \textbf{Stacked architectures:} multiple FlowNet modules chained (FlowNetC + FlowNetS)
\item \textbf{Small displacement module (FlowNet-SD):} specialized sub-network for slow motions
\item \textbf{Multi-stage training schedule} on multiple datasets
\end{itemize}
\begin{center}\includegraphics[width=0.600\linewidth]{Pasted_image_20260522135728.png}\end{center}
\begin{center}\includegraphics[width=0.600\linewidth]{Pasted_image_20260522135809.png}\end{center}

\begin{center}
\begin{tabular}{|l|l|}
\hline
\textbf{Aspect} & \textbf{Detail} \\ \hline
Quality & Significantly lower EPE than FlowNet \\ \hline
Size & Large memory footprint (stacked networks) \\ \hline
Efficiency & Trade-off: FlowNet2 (accurate) vs FlowNet2-s (faster, lighter) \\ \hline
\end{tabular}
\end{center}

\subsection{PWC-Net (2018)}

\textgreater{} Sun et al. (2018). \textit{PWC-Net: CNNs for Optical Flow Using Pyramid, Warping, and Cost Volume}

\textbf{Hybrid design} — embeds classical vision priors into CNN:

Three innovations:
\begin{itemize}
\item \textbf{P}yramidal feature extraction (coarse-to-fine)
\item \textbf{W}arping: warp reference features using current flow estimate
\item \textbf{C}ost volume: local cross-correlation between warped and current features (limited search range)
\end{itemize}

\begin{center}
\begin{tabular}{|l|l|}
\hline
\textbf{Aspect} & \textbf{Detail} \\ \hline
Parameters & 17× fewer than FlowNet2 \\ \hline
Speed & ≈ 30ms for Sintel-resolution images \\ \hline
Evaluation & SOTA on KITTI benchmarks; strong sim-to-real generalization \\ \hline
\end{tabular}
\end{center}

\subsection{RAFT (2020)}

\textgreater{} Teed \& Deng (2020). \textit{RAFT: Recurrent All-Pairs Field Transforms for Optical Flow}

\textbf{Innovations:}

\begin{itemize}
\item \textbf{All-pairs correlation volume:} compute correlation between \textit{every pair} of pixels (4D tensor)
\item \textbf{GRU-based iterative update:} recurrent refinement of a single high-resolution flow field
\item \textbf{No spatial pyramids:} avoids detail loss from coarse-to-fine
\end{itemize}

\begin{calloutbox}{red!70!black}{RAFT Key Property}
Maintains full-resolution flow field throughout — no pyramid-induced information loss. Number of GRU iterations tunable at inference (speed-quality trade-off).
\end{calloutbox}

\begin{center}
\begin{tabular}{|l|l|}
\hline
\textbf{Aspect} & \textbf{Detail} \\ \hline
Accuracy & SOTA on Sintel and KITTI; precise on thin structures and boundaries \\ \hline
Generalization & Unprecedented "zero-shot" transfer from synthetic to real video \\ \hline
Speed & Slower than PWC-Net \\ \hline
Complexity & All-pairs correlation is memory-intensive (mitigated via Correlation Pyramid) \\ \hline
\end{tabular}
\end{center}

\subsection{DL Perspectives and Challenges}

\textbf{Main advantages:}
\begin{itemize}
\item Robust to occlusions and large displacements
\item High precision in textureless areas
\end{itemize}

\textbf{Current challenges:}

\begin{center}
\begin{tabular}{|l|l|}
\hline
\textbf{Challenge} & \textbf{Description} \\ \hline
Generalization & Performance drops on data outside training distribution \\ \hline
Computational cost & High memory/power for SOTA models \\ \hline
Integration & Real-time and video coding standards require accuracy-efficiency balance \\ \hline
\end{tabular}
\end{center}

\textit{DL dominates analysis tasks (optical flow); block-matching remains standard for video compression.}

\subsection{Evolutionary Summary}

\begin{center}
\begin{tabular}{|l|l|l|}
\hline
\textbf{Model} & \textbf{Year} & \textbf{Key Innovation} \\ \hline
FlowNet & 2015 & First end-to-end CNN for dense flow; correlation layer \\ \hline
FlowNet 2.0 & 2017 & Stacked networks; small displacement module \\ \hline
PWC-Net & 2018 & Pyramids + Warping + Cost Volume; 17× smaller \\ \hline
RAFT & 2020 & All-pairs correlation + GRU iterative update; no pyramids \\ \hline
\end{tabular}
\end{center}

\bigskip\hrule\bigskip

\section{6. Conclusions}

\begin{itemize}
\item \textbf{Motion Estimation} extracts motion information from video sequences
\item Focus is on \textbf{optical flow} (apparent motion), not physical motion
\item \textbf{Block-matching:} conceptually simple, widely used in video compression standards
\item \textbf{Neural networks:} emerging paradigm, dominant for analysis tasks
\end{itemize}
\begin{center}\includegraphics[width=0.600\linewidth]{Pasted_image_20260522135856.png}\end{center}
\bigskip\hrule\bigskip

\section{Summary Table}

\begin{center}
\begin{tabular}{|l|l|l|l|l|l|}
\hline
\textbf{Method} & \textbf{Type} & \textbf{Motion Model} & \textbf{Output} & \textbf{Key Advantage} & \textbf{Key Limitation} \\ \hline
Horn-Schunck & Variational & Continuous smooth field & Dense (1 vec/pixel) & Convex, unique solution & Over-smooth at boundaries \\ \hline
Block Matching (FS) & Block-based & Translation per block & Sparse (1 vec/block) & Global optimum guaranteed & $O(n^2)$ complexity \\ \hline
Block Matching (fast) & Block-based & Translation per block & Sparse & Low complexity & Local minima risk \\ \hline
Affine Parametric & Parametric & Affine (6 params) & Global/regional & Very compact representation & Requires segmentation or region \\ \hline
FlowNet & DL & Learned end-to-end & Dense & Fast inference (GPU) & Weak on small displacements \\ \hline
PWC-Net & DL (hybrid) & Pyramidal + Cost Volume & Dense & 17× smaller than FlowNet2; SOTA KITTI & Loss at pyramid boundaries \\ \hline
RAFT & DL (recurrent) & All-pairs + GRU & Dense & Best accuracy; zero-shot generalization & Memory-intensive \\ \hline
\end{tabular}
\end{center}

\begin{center}
\begin{tabular}{|l|l|}
\hline
\textbf{Parameter} & \textbf{Trade-off} \\ \hline
Block size $P \times Q$ ↑ & Complexity ↓, Coding cost ↓, MSE ↑ \\ \hline
Search window $\mathcal{W}$ ↑ & Accuracy ↑, Complexity ↑ \\ \hline
Subpixel precision ↑ & Accuracy ↑, Complexity ↑ \\ \hline
$p = 1$ (SAD) & Regular MVF, lower coding cost, worse MSE \\ \hline
$p = 2$ (SSD) & Better MSE, irregular MVF, higher coding cost \\ \hline
$\lambda_{ME}$ ↑ & Coding cost ↓, MSE ↑ \\ \hline
\end{tabular}
\end{center}

